\paragraph{Order \(\LUV^2\): local frame rotations, the Goldstone-extremised density and the Noether coefficient of \(\rho^2X\) (conditional; not closure).}
Keep the setting of the preceding paragraph and of the paragraph on the torsion invariant: Poisson limit, the \(\LUV^2\) term, \(e^A{}_\mu=\eta^{AB}\omega_{\mu B}\) with \(\omega=\Theta^{-1}\) closed, \(V_\mu=G^{\lambda\nu}T_{\lambda\mu\nu}\), equal to \(\partial_\mu\ln\rho\) on the Poisson class, \(\mathcal T=\tfrac14T_{\lambda\mu\nu}T^{\lambda\mu\nu}+\tfrac12T_{\lambda\mu\nu}T^{\mu\lambda\nu}-V_\mu V^\mu\), the density \(\tfrac{N_{\mathrm{tot}}}6\rho R[g_\rho]+\kappa_X\rho^2X\), and the weights \(\kappa_X=1-n_D/2\).

\emph{Classification.} Under \(\delta e^A=\lambda^A{}_B(x)e^B\), \(\lambda_{AB}=-\lambda_{BA}\), the metric \(G\) and \(\rho\), with \(\rho^2=\lvert\det e\rvert\) for a general coframe, are unchanged and \(\sqrt{\lvert G\rvert}\,(a_1T_{\lambda\mu\nu}T^{\lambda\mu\nu}+a_2T_{\lambda\mu\nu}T^{\mu\lambda\nu}+a_3V_\mu V^\mu+a_4V^\mu\partial_\mu\ln\rho+a_5\lvert\partial\ln\rho\rvert^2)\) of an arbitrary coframe changes by \(\sqrt{\lvert G\rvert}\,P^{\mu|\lambda\nu}\partial_\mu\lambda_{\lambda\nu}\), with \(P\) linear in \(T\) (checked against a finite variation). The change is a total derivative for every coframe only on the span of \(\mathcal T\) and \(\lvert\partial\ln\rho\rvert^2\), that is \((a_1,a_2,a_3,a_4)\propto(\tfrac14,\tfrac12,-1,0)\) with \(a_5\) free: the null space of the divergence of \(\sqrt{\lvert G\rvert}P\) at random two-jets has dimension two and is this span, for \(d=3,4\) and both signatures, and at quadratic order about a constant frame the invariance conditions have the exact rational solution \(a_1=-a_3/4\), \(a_2=-a_3/2\). Both members are functions of \(G\) up to total derivatives. The exact algebraic identity \(X=2\mathcal T+\tfrac32S_{\lambda\mu\nu}S^{\lambda\mu\nu}+2V_\mu V^\mu\) of the irreducible torsion parts then shows that, with \(V_\mu V^\mu\) replaced by its class value \(\lvert\partial\ln\rho\rvert^2\), the non-invariance of \(\kappa_X\rho^2X\) is the axial term \(\tfrac32\kappa_X\rho^2S^2\) (the difference \(2(V_\mu V^\mu-\lvert\partial\ln\rho\rvert^2)\) vanishes on the class but is not invariant); on generic Poisson-class two-jets its rotation variation is nonzero, and no addition of multiples of \(V_\mu-\partial_\mu\ln\rho\), which vanishes on the class, reaches the invariant span. The frame sector of the induced density is therefore not invariant under local rotations of the \(\Theta\)-built frame unless \(\kappa_X=0\), that is \(n_D=2\) in the weights of the paragraph on fermionic content.

\emph{Dressed density.} Let the Stueckelberg field \(\Lambda\) of Equation~\eqref{eq:dressed-omega} act on the coframe, \(e\to\Lambda e\), with the density \(\kappa_X\rho^2X[\Lambda e]\); this is an input of Assumption~\ref{ass:local-lorentz} and is not derived. The density is then invariant under \(e\to\Lambda'e\), \(\Lambda\to\Lambda\Lambda'^{-1}\) of the pair, and \(\Lambda\) is treated as a field fixed by its own stationarity at the \(\Theta\)-built frame \(e\). In the unitary gauge \(\Lambda=\mathbb I\) with \(e\) that frame the density is the undressed one and \(r\) stays undetermined, as in the preceding paragraph. About constant \(\bar\Theta\), for plane-wave perturbations of the Poisson class (random constant \(\bar\Theta\) and wave vectors, both signatures), the stationary point has \(S=0\), leaves \(\mathcal T\) and \(\lvert\partial\ln\rho\rvert^2\) unchanged, and replaces \(\rho^2X\) by \(2\mathcal T+2V_\mu V^\mu\): one half of the on-class value on the two modes with \(\delta\rho=0\) (ratio \(0.5\) to \(10^{-12}\)). The reason is that \(S\) is annihilated by the wedge with the wave vector while the rotations shift it by the full image of that wedge. The Hessian over the twelve real rotation amplitudes is positive definite for \(\eta=\delta\) and indefinite, with six positive and six negative eigenvalues among twelve real parameters, in the Lorentzian cases tested. On the family \(\Theta^{12}=1/f\), \(\Theta^{34}=1\) the axial torsion vanishes pointwise, \(X=2\lvert\partial\ln\rho\rvert^2_G\), and the rotation equation holds identically (\(n=0.8,1.2\)), so the family lies on the stationary set without extremisation.

\emph{Noether coefficient.} At the stationary points above, where \(S=0\) and \(V_\mu=\partial_\mu\ln\rho\), the density depends on \(e\) only through \(G\) and \(\rho\), and \(\mathcal T\simeq-R\) up to a total derivative, so \(\rho^2X\to-2R+2V_\mu V^\mu\): the joint flow whose internal part is a local symmetry selects \(r=-2\), and \(c_W=N_{\mathrm{tot}}/6-2\kappa_X\) equals the tensor coefficient \(c_T^{\mathrm{ext}}=N_{\mathrm{tot}}/6-2\kappa_X\) of the dressed form. The value \(r=-4\) is the on-class tensor coefficient with \(S^2\) retained, for which the compensating rotation is not a symmetry and the construction of the preceding paragraph is not available; any value of \(r\) other than \(-2\) corresponds to a different choice of the non-invariant remainder, a choice that the symmetry does not make. Hence \(c_A-c_T^{\mathrm{ext}}=2\kappa_X\) for the fixed-state replica at \(n_D=0\), and \(c_A-c_T^{\mathrm{ext}}=2\) for the antiperiodic Dirac gluing at every \(N_s\) and \(n_D\), smaller by \(2\kappa_X\) than the value \(4-n_D\) before the dressing, that is half of it at \(n_D=0\) and unchanged at \(n_D=2\); with the vector holonomy contact \(-\alpha_v\sigma\), \(\alpha_v=1\), the residue is \(2-\alpha_v\sigma\), that is \(1\) on \(n<1\) and \(3\) on \(n>1\). A contact of modulus \(2\) would close it; no such contact is derived here.

These statements are for the Poisson-limit density, the \(\LUV^2\) term and the weights above, quadratic order about constant \(\bar\Theta\) for the plane-wave class, the stationary set of the rotation parameters (tree level), and the one family named; they do not derive the Stueckelberg field or the local rotation symmetry from the matrix model, the independent diagonal connection \(B\) or a covariant bundle connection for the fluctuation operator, the stationary set beyond the cone family, the loop-level treatment of \(\Lambda\), the entropy of the rigid compensating rotation, the replica or the holonomy selected by the matrix state, or a value or sign of \(G_{\mathrm{eff}}\). Problems~B and~D remain open and PROGRAMME-CONDITIONAL.
